\section{可持久化线段树}

\subsection{可持久化数组（P3919）}

\par \noindent 用主席树维护每个版本的数组。每次修改从旧版本派生新版本，只复制被修改路径上的节点，其余节点与原版本共用。

\begin{minted}{c++}
const int N=3e7+10;
struct seg
{
    int l,r,v;
}t[N];
int a[N];
int now;

void build(int &p,int l,int r)
{
    p=++now;
    if(l==r)
    {
        t[p].v=a[l];
        return;
    }
    int mid=(l+r)>>1;
    build(t[p].l,l,mid),build(t[p].r,mid+1,r);
}
void modify(int &p,int l,int r,int x,int y)
{
    t[++now]=t[p];p=now;
    if(l==r)
    {
        t[p].v=y;
        return;
    }
    int mid=(l+r)>>1;
    if(x<=mid)modify(t[p].l,l,mid,x,y);
    if(x>mid)modify(t[p].r,mid+1,r,x,y);
}
int query(int p,int l,int r,int x)
{
    if(l==r)return t[p].v;
    int mid=(l+r)>>1;
    if(x<=mid)return query(t[p].l,l,mid,x);
    if(x>mid)return query(t[p].r,mid+1,r,x);
}
int rt[N];

signed main()
{
    int n=read(),m=read();
    for(int i=1;i<=n;i++)a[i]=read();
    build(rt[0],1,n);
    for(int i=1;i<=m;i++)
    {
        int his=read(),op=read();
        if(op==1)
        {
            int x=read(),y=read();
            rt[i]=rt[his];
            modify(rt[i],1,n,x,y);
        }
        else
        {
            int x=read();
            cout<<query(rt[his],1,n,x)<<endl;
            rt[i]=rt[his];
        }
    }
    return 0;
}
\end{minted}

\subsection{静态区间第 k 小（P3834）}

\par \noindent 对序列建主席树，第 $i$ 个版本维护前 $i$ 个数的权值线段树。查询区间 $[l,r]$ 第 $k$ 小时，用版本 $r$ 减去版本 $l-1$ 得到区间内的权值分布，再在权值线段树上二分。

\begin{minted}{c++}
const int N=3e7+10;
struct seg
{
    int l,r,v;
}t[N];
int a[N],b[N];
int now;

void build(int &p,int l,int r)
{
    p=++now;
    if(l==r)return;
    int mid=(l+r)>>1;
    build(t[p].l,l,mid),build(t[p].r,mid+1,r);
}
void modify(int &p,int l,int r,int x)
{
    t[++now]=t[p];p=now;t[p].v++;
    if(l==r)return; 
    int mid=(l+r)>>1;
    if(x<=mid)modify(t[p].l,l,mid,x);
    if(x>mid)modify(t[p].r,mid+1,r,x);
}
int query(int p1,int p2,int l,int r,int k)
{
    if(l==r)return b[l];
    int s=t[t[p2].l].v-t[t[p1].l].v;
    int mid=(l+r)>>1;
    if(s>=k)return query(t[p1].l,t[p2].l,l,mid,k);
    if(s<k)return query(t[p1].r,t[p2].r,mid+1,r,k-s);
}
int rt[N];

signed main()
{
    int n=read(),q=read();
    for(int i=1;i<=n;i++)a[i]=read(),b[i]=a[i];
    sort(b+1,b+1+n);int m=unique(b+1,b+1+n)-b-1;
    for(int i=1;i<=n;i++)a[i]=lower_bound(b+1,b+1+m,a[i])-b;
    build(rt[0],1,m);
    for(int i=1;i<=n;i++)
    {
        rt[i]=rt[i-1];
        modify(rt[i],1,m,a[i]);
    }
    for(int i=1;i<=q;i++)
    {
        int l=read(),r=read(),k=read();
        cout<<query(rt[l-1],rt[r],1,m,k)<<endl;
    }
    return 0;
}
\end{minted}

\section{可持久化并查集}
\begin{tcolorbox}
给定 $n$ 个集合，第 $i$ 个集合内初始状态下只有一个数，为 $i$。
有 $m$ 次操作。操作分为 3 种：
\begin{itemize}
\item 1 $a$ $b$ 合并 $a,b$ 所在集合
\item 2 $k$ 回到第 $k$ 次操作（执行三种操作中的任意一种都记为一次操作）之后的状态
\item 3 $a$ $b$ 询问 $a,b$ 是否属于同一集合，如果是则输出 1，否则输出 0。
\end{itemize}
\end{tcolorbox}

\par \noindent 可持久化并查集不支持路径压缩，只能按秩合并。用可持久化数组维护每个版本的 $fa$ 和 $sz$，每次操作从上一个版本派生新版本。

\begin{minted}{c++}
using namespace std;
using pii=pair<int,int>;
const int N=200005;
int n,m;

struct Array
{
    struct Segment
    {
        int l,r,v;
    }t[N*80];
    int rt[N],cnt;
    void build(int &u,int l,int r,int c)
    {
        t[u=++cnt]={0,0,0};
        if(l==r) 
        {
            if(c==1) t[u].v=l;
            else if(c==2) t[u].v=1;
            return;
        }
        int mid=l+r>>1;
        build(t[u].l,l,mid,c),build(t[u].r,mid+1,r,c);
    }
    void update(int &u,int pre,int l,int r,int pos,int v)
    {
        t[u=++cnt]=t[pre];
        if(l==r) return t[u].v=v,void();
        int mid=l+r>>1;
        if(pos<=mid) 
            update(t[u].l,t[pre].l,l,mid,pos,v);
        else
            update(t[u].r,t[pre].r,mid+1,r,pos,v);
    }
    int query(int u,int l,int r,int pos)
    {
        if(l==r) return t[u].v;
        int mid=l+r>>1;
        if(pos<=mid) 
            return query(t[u].l,l,mid,pos);
        else
            return query(t[u].r,mid+1,r,pos);
    }
    void clear(){cnt=0;}
    int inline get(int u,int pos)
    {
        return query(rt[u],1,n,pos);
    }
    void inline ins(int u,int pos,int v)
    {
        update(rt[u],rt[u],1,n,pos,v);
    }
}sz,fa;

int find(int u,int x)
{
    int f=fa.get(u,x);
    return x==f?x:find(u,f);
}
bool inline merge(int v,int a,int b)
{
    a=find(v,a),b=find(v,b);
    if(a==b) return false;
    int sa=sz.get(v,a),sb=sz.get(v,b);
    if(sa>sb) swap(a,b),swap(sa,sb);
    fa.ins(v,a,b);sz.ins(v,b,sa+sb);
    return true;
}

signed main()
{
    ios::sync_with_stdio(false);cin.tie(nullptr);cout.tie(nullptr);

    cin>>n>>m;
    fa.build(fa.rt[0],1,n,1);
    sz.build(sz.rt[0],1,n,2);

    int op,x,y;
    for(int i=1;i<=m;i++) 
    {
        cin>>op>>x;
        if(op==1)
        {
            cin>>y;
            fa.rt[i]=fa.rt[i-1];
            sz.rt[i]=sz.rt[i-1];
            merge(i,x,y);
        }
        else if(op==2)
        {
            fa.rt[i]=fa.rt[x];
            sz.rt[i]=sz.rt[x];
        }
        else 
        {
            cin>>y;
            int fx=find(i-1,x),fy=find(i-1,y);
            cout<<int(fx==fy)<<'\n';
            fa.rt[i]=fa.rt[i-1];
            sz.rt[i]=sz.rt[i-1];
        }
    }
    return 0;
}
\end{minted}

\par \noindent 说明：
\begin{itemize}
\item \textbf{可持久化数组} \texttt{Array}：用主席树维护每个版本的数组。\texttt{build} 初始化第 0 个版本，\texttt{update} 从旧版本派生新版本，\texttt{query} 查询某个版本的某个位置。
\item \textbf{两个数组}：\texttt{fa} 维护每个点的父亲，\texttt{sz} 维护每个集合的大小。
\item \textbf{find}：沿 \texttt{fa} 向上找根。因为不能路径压缩，所以递归查询 \texttt{fa.get(u,x)}。
\item \textbf{merge}：按大小合并，把小的接到大的下面，并更新 \texttt{sz}。
\item \textbf{版本派生}：每次操作先把当前版本的 \texttt{rt} 复制成上一版本的 \texttt{rt}，再在这个基础上修改，得到新版本。
\item \textbf{操作 2}：直接令第 $i$ 个版本的 \texttt{rt} 指向第 $x$ 个版本的 \texttt{rt}，实现回退。
\end{itemize}